\chapter{等离子体效应}
\
本节主要讨论电磁波在冷等离子体中的传播效应。
所谓的“冷”等离子体，是指等离子体粒子的热运动速度远小于波的相速度，且介质自身的辐射和吸收相对于入射场可忽略，主要关注波的传播特性（如色散、折射、法拉第旋转等）。
\n{好比正弦函数与一个幅度远低于它的正弦函数叠加，幅值不太变但是相位改变}
没有散射和吸收，冷等离子气体与入射光子场的相互作用与之前有很大不同。之前的辐射是对运动电荷辐射的计算（如韧致辐射、同步辐射），而现在是考察冷等离子体在入射场作用下所导致的电荷与电流对传播的影响，其中等离子体的发射和吸收相对入射场是可以忽略的，电子运动只是一种次级效应。

\section{各向同性冷等离子体中的色散关系}
\n{金属就是致密的等离子气体。}
在无外部磁场的情况下，等离子体是各向同性的，这里把冷等离子体受入射场影响产生的次级效应归于介电常数中去。

\subsection{介电常数与等离子体频率}
\paragraph{冷等离子体的电导率}
\n{假设离子质量远大于电子，忽略离子的运动，仅考虑电子对电流的贡献。}
首先考察单个电子在时谐电磁场 $\vb{E}(\vb{r}, t) = \vb{E}_0 e^{i(\vb{k}\cdot\vb{r} - \omega t)}$ 作用下的运动。
忽略非相对论情况下的洛伦兹力磁场项 ($v/c \ll 1$)，运动方程为：
\begin{equation}
    m \frac{d\vb{v}}{dt} = -e \vb{E} \implies -i\omega m \vb{v} = -e \vb{E}
\end{equation}
解得电子速度 $\vb{v} = \frac{e}{i\omega m}\vb{E}$。由此得到电流密度 $\vb{j}$ 与电导率 $\sigma$ 的关系：
\begin{equation}
    \vb{j} = -n e \vb{v} = \frac{i n e^2}{\omega m} \vb{E} \equiv \sigma \vb{E}
\end{equation}
其中电导率 $\sigma = \frac{i n e^2}{\omega m}$ 为纯虚数。

\paragraph{介电常数}

\n{在傅里叶空间中，$\nabla \to i\vb{k}$，$\partial_t \to -i\omega$。}
考虑源项为感生电流 $\vb{j} = \sigma \vb{E}$ 的麦克斯韦方程组。在 $(k, \omega)$ 空间中，两个旋度方程写为：
\begin{align}
    i\vb{k} \times \vb{E} &= \frac{i\omega}{c} \vb{B} \label{eq:maxwell_k_E} \\
    i\vb{k} \times \vb{B} &= \frac{4\pi}{c} \vb{j} - \frac{i\omega}{c} \vb{E} = \frac{4\pi \sigma}{c} \vb{E} - \frac{i\omega}{c} \vb{E} \label{eq:maxwell_k_B}
\end{align}

为了得到关于 $\vb{E}$ 的波动方程，我们将 \eqref{eq:maxwell_k_E} 式中的 $\vb{B}$ 代入 \eqref{eq:maxwell_k_B} 式：
\begin{equation}
    i\vb{k} \times \left( \frac{c}{\omega} \vb{k} \times \vb{E} \right) = - \frac{i\omega}{c} \left( 1 + \frac{4\pi i \sigma}{\omega} \right) \vb{E}
\end{equation}
整理上式，引入\textbf{介电常数} $\epsilon(\omega)$：
\begin{equation}
    \vb{k} \times (\vb{k} \times \vb{E}) = - \frac{\omega^2}{c^2} \epsilon(\omega) \vb{E}
\end{equation}
其中介电常数定义为：
\begin{equation}
    \epsilon(\omega) \equiv 1 + \frac{4\pi i \sigma}{\omega} = 1 - \frac{4\pi n e^2}{m \omega^2} = 1 - \frac{\omega_p^2}{\omega^2}
\end{equation}
此处 $\omega_p = \sqrt{4\pi n e^2 / m}$ 为等离子体频率。
\paragraph{色散关系}
应用横波条件 $\vb{k} \cdot \vb{E} = 0$，波动方程简化为：
\begin{equation}
    -k^2 \vb{E} = - \frac{\omega^2}{c^2} \epsilon(\omega) \vb{E}
\end{equation}
消去 $\vb{E}$，即得到最终的\textbf{色散关系}：
\begin{equation}
    k^2 = \frac{\omega^2}{c^2} \epsilon(\omega) = \frac{\omega^2}{c^2} \left( 1 - \frac{\omega_p^2}{\omega^2} \right)
\end{equation}
或者写作：
\begin{equation}
    \omega^2 = \omega_p^2 + c^2 k^2
\end{equation}
在不同的频率下，电磁波传播有两种模式：
    \begin{itemize}
        \item \textbf{传播模式 ($\omega > \omega_p$)}: $k$ 为实数，波可以在介质中传播。此时相速度 $v_{ph} = \omega/k = c/\sqrt{\epsilon} > c$，群速度 $v_g = d\omega/dk = c\sqrt{\epsilon} < c$。
        \item \textbf{截止模式 ($\omega < \omega_p$)}: $\epsilon < 0$，$k$ 为纯虚数。波幅随距离指数衰减（趋肤效应），能量无法传入介质深处，发生全反射。
        \n{电离层反射无线电波即基于此原理。}
    \end{itemize}

\paragraph{相速度与群速度}
色散关系导致了波的相速度与群速度分离。首先定义等离子体的\textbf{折射率} $n_r$ 为介电常数的平方根：
\begin{equation}
    n_r(\omega) \equiv \sqrt{\epsilon(\omega)} = \sqrt{1 - \frac{\omega_p^2}{\omega^2}}
\end{equation}
\textbf{相速度}描述单色波等相面的推进速度。由于在传播模式下 $n_r < 1$，相速度表现出“超光速”特性：
\begin{equation}
    v_{ph} \equiv \frac{\omega}{k} = \frac{c}{n_r} > c
\end{equation}
这并不违反相对论，因为纯单色波不携带起始信息。能量与信号的传播速度由\textbf{群速度}描述，通过对色散关系 $\omega^2 = \omega_p^2 + c^2k^2$ 微分可得：
\begin{equation}
    v_g \equiv \frac{d\omega}{dk} = c^2 \frac{k}{\omega} = c n_r
\end{equation}
由此可见，在冷等离子体中，相速度与群速度始终满足以下互补关系：
\begin{equation}
    v_{ph} v_g = c^2
\end{equation}
当入射频率接近等离子体频率（$\omega \to \omega_p$）时，$n_r \to 0$，导致 $v_{ph} \to \infty$ 而 $v_g \to 0$，波包能量停止传播并发生全反射。
\section{脉冲星色散测量}

由于群速度 $v_g(\omega)$ 依赖于频率，宽带脉冲信号（如脉冲星辐射）在星际介质中传播时，高频成分比低频成分先到达地球。

到达时间 $t_p$ 为路径积分：
\begin{equation}
    t_p(\omega) = \int_0^d \frac{ds}{v_g} = \frac{1}{c} \int_0^d \left( 1 - \frac{\omega_p^2}{\omega^2} \right)^{-1/2} ds
\end{equation}
在射电波段通常满足 $\omega \gg \omega_p$，利用泰勒展开 $(1-x)^{-1/2} \approx 1 + x/2$：
\begin{equation}
    t_p \approx \frac{d}{c} + \frac{1}{2c\omega^2} \int_0^d \omega_p^2 ds = \frac{d}{c} + \frac{2\pi e^2}{mc\omega^2} \underbrace{\int_0^d n_e ds}_{\text{DM}}
\end{equation}
定义\textbf{色散量}为视线方向上的电子柱密度：
\begin{equation}
    \text{DM} \equiv \int_0^d n_e ds
\end{equation}
通过测量不同频率到达的时间延迟 $\Delta t$，可以求得 DM，进而估算脉冲星距离。

\section{折射介质中的辐射转移}

在非均匀介质（如星际云或恒星大气）中，折射率 $n(\vb{r})$ 随空间位置变化，光线不再沿直线传播，而是发生折射（弯曲）。
此时，即使介质没有吸收和发射，辐射强度 $I_\nu$ 也不再是沿光线的守恒量（因为光束的截面和张角发生了拉伸或压缩）。
我们需要寻找一个新的绝热不变量来建立转移方程。

\begin{theorem}
    光线穿过介质时，
    \begin{align*}
        \frac{I_\nu}{n^2} = const
    \end{align*}
    是绝热不变量。
\end{theorem}
\begin{proof}
    考虑辐射能量 $dE$ 穿过两个折射率分别为 $n_1$ 和 $n_2$ 的介质分界面。
设光线入射角为 $\theta_1$，折射角为 $\theta_2$。根据斯涅尔定律：
\begin{equation}
    n_1 \sin\theta_1 = n_2 \sin\theta_2
\end{equation}
对上式微分，并结合方位角 $d\phi$ 不变，由 $d\Omega = \sin\theta d\theta d\phi$ 可得：
\[ (n_1 \sin\theta_1)(n_1 \cos\theta_1 d\theta_1) d\phi = (n_2 \sin\theta_2)(n_2 \cos\theta_2 d\theta_2) d\phi \]
    利用 $d\Omega = \sin\theta d\theta d\phi$，整理得：
    \begin{equation}
        n_1^2 \cos\theta_1 d\Omega_1 = n_2^2 \cos\theta_2 d\Omega_2
    \end{equation}
    根据能量守恒，通过界面的能量流必须相等：
\begin{equation}
    dE_1 = dE_2 \implies I_{\nu,1} \cos\theta_1 d\Omega_1 dA = I_{\nu,2} \cos\theta_2 d\Omega_2 dA
\end{equation}
将立体角关系代入能量方程，消去公因子 $dA \cos\theta d\Omega$，得到：
\begin{equation}
    \frac{I_{\nu,1}}{n_1^2} = \frac{I_{\nu,2}}{n_2^2} \implies \frac{I_\nu}{n_\nu^2} = \text{const}
    \label{eq:In}
\end{equation}
\end{proof}

\subsection{考虑折射率的辐射转移方程}
对\ref{eq:In}求偏微分得：
\begin{align*}
    \frac{I_\nu}{\partial s} = \frac{2 I_\nu}{n_\nu}\frac{\partial n_\nu}{\partial s}
\end{align*}
这就是说，将这个关于折射率得偏导项加入辐射转移方程可得：
\begin{equation}
    \dv{I_\nu}{s} = j_\nu - \alpha_\nu I_\nu + \underbrace{\frac{I_\nu}{n_\nu^2} \dv{n_\nu^2}{s}}_{\text{折射发散项}}
\end{equation}
\n{可以看出考虑折射率变化后，方程就是把原来的强度量纲的量除以$\nu^2$，如$I_\nu$替换为$I_\nu/\nu^2$，}
简化为：
\begin{equation}
    n_\nu^2 \dv{s} \left( \frac{I_\nu}{n_\nu^2} \right) = j_\nu - \alpha_\nu I_\nu
\end{equation}
定义源函数为，可简化辐射转移方程：
\begin{align}
    S_\nu \equiv \frac{1}{n_{\nu}^2} \frac{j_\nu}{\alpha_\nu} \Rightarrow \frac{d}{d\tau_\nu}\left(\frac{I_\nu}{n_\nu^2}\right) = -\frac{I_\nu}{n^2_\nu} +S_\nu
\end{align}

\section{磁化冷等离子体：法拉第旋转}

当等离子体中存在大尺度背景磁场 $\vb{B}_0$ 时，空间的对称性被破坏，介质变为\textbf{各向异性}。洛伦兹力项 $-e(\vb{v} \times \vb{B}_0)/c$ 不再能被忽略，这导致介质对左旋和右旋偏振光的响应不同，进而产生法拉第旋转效应。

\n{注意：这里假设离子静止，仅考虑电子运动。}

\subsection{光线在磁化冷等离子体中的传播}

当等离子体中存在强背景磁场 $\vb{B}_0$ 时，介质呈现各向异性。我们需要从单电子的动力学出发，推导其传播特性。

\subsubsection{电子运动方程与圆偏振分解}

考虑一个沿磁场方向传播的电磁波（$\vb{k} \parallel \vb{B}_0$）。设定 $\vb{B}_0 = B_0 \hat{z}$，$\vb{k} = k \hat{z}$。
电子在电磁波场中的非相对论运动方程为：
\begin{equation}
    m \dot{\vb{v}} = -e \vb{E} - \frac{e}{c} \vb{v} \times \vb{B}_0
\end{equation}

\n{此处作如下近似：1. 非相对论 ($v \ll c$)；2. 忽略波本身的磁场 $\vb{B}_{wave}$，假设背景磁场 $B_0 \gg B_{wave}$，即线性化处理。}

在各向异性介质中，$\vb{v}$ 与 $\vb{E}$ 的方向不再单纯平行，简单的时谐假设 $\vb{E} \propto e^{-i\omega t}$ 会导致复杂的张量运算。
物理上，由于电子绕磁场线做圆周运动，最自然的基底不是直角坐标 $(x, y)$，而是\textbf{圆偏振坐标}。

我们假设电磁波是圆偏振的，将矢量分解为右旋 ($R$) 和左旋 ($L$) 分量。定义复矢量 $v_\pm = v_x \pm i v_y$ 和 $E_\pm = E_x \pm i E_y$（假设时变因子为 $e^{-i\omega t}$）。
将运动方程在 $x, y$ 方向展开并重新组合：
\begin{align}
    -i\omega m v_x &= -e E_x - \frac{e B_0}{c} v_y \\
    -i\omega m v_y &= -e E_y + \frac{e B_0}{c} v_x
\end{align}
将第二式乘以 $\pm i$ 加到第一式上，利用 $\omega_B = eB_0/mc$，可得解耦后的标量方程：
\begin{equation}
    -i\omega m (v_x \pm i v_y) = -e (E_x \pm i E_y) \mp i \frac{e B_0}{c} (v_x \pm i v_y)
\end{equation}
整理得到电子速度的解：
\begin{equation}
    v_\pm = \frac{e E_\pm}{i m (\omega \mp \omega_B)}
\end{equation}

\subsubsection{介电常数与色散关系}

利用电流密度 $j_\pm = -n e v_\pm = \sigma_\pm E_\pm$，可得电导率 $\sigma_\pm$：
\begin{equation}
    \sigma_\pm = \frac{i n e^2}{m(\omega \mp \omega_B)}
\end{equation}
代入介电常数定义 $\epsilon = 1 + 4\pi i \sigma / \omega$，得到两种模式的\textbf{介电常数}：
\begin{equation}
    \epsilon_R(\omega) = 1 - \frac{\omega_p^2}{\omega(\omega - \omega_B)}, \quad
    \epsilon_L(\omega) = 1 - \frac{\omega_p^2}{\omega(\omega + \omega_B)}
\end{equation}

由此得到沿磁场传播的\textbf{色散关系}（Dispersion Relation）：
\begin{equation}
    k_{R,L}^2 = \frac{\omega^2}{c^2} \epsilon_{R,L}(\omega) = \frac{\omega^2}{c^2} \left[ 1 - \frac{\omega_p^2}{\omega(\omega \mp \omega_B)} \right]
\end{equation}

\subsubsection{相速度与群速度}

由于 $\epsilon_R \neq \epsilon_L$，两种偏振模具有不同的传播速度。

\begin{equation}
    v_{ph, R/L} = \frac{\omega}{k_{R/L}} = \frac{c}{\sqrt{\epsilon_{R,L}}} = c \left[ 1 - \frac{\omega_p^2}{\omega(\omega \mp \omega_B)} \right]^{-1/2}
\end{equation}
\n{R 模在 $\omega = \omega_B$ 处发生共振，$\epsilon_R \to \infty$，意味着波速趋于 0，能量被剧烈吸收。}

群速度 $v_g = d\omega / dk = (dk/d\omega)^{-1}$。对色散关系微分计算较为繁琐，通式为：
\begin{equation}
    v_{g, R/L} = \frac{c}{\frac{d}{d\omega}(\omega \sqrt{\epsilon_{R,L}})} = \frac{c}{\sqrt{\epsilon_{R,L}} + \frac{\omega}{2\sqrt{\epsilon_{R,L}}} \frac{d\epsilon_{R,L}}{d\omega}}
\end{equation}
在远离回旋共振频率（$\omega \gg \omega_B$）的极限下，群速度不仅依赖于频率，还受到磁场的强烈调制。这种群速度的差异是导致脉冲信号在磁化等离子体中发生畸变的根本原因。

\subsection{法拉第旋转 (Faraday Rotation)}

一束入射的\textbf{线偏振光}可以等效分解为一束左旋波和一束右旋波的叠加，然而经过磁化冷等离子气体时，两部分的相速度不同，那么出射的时候合成的线偏振光就会出现一定程度的旋转。

旋转角 $\Delta \theta$ 等于两个圆偏振分量相位差的一半：
\begin{equation}
    \Delta \theta = \frac{1}{2} (\phi_R - \phi_L) = \frac{1}{2} \int_0^d (k_R - k_L) ds
\end{equation}
在典型的天体物理环境中（射电波段），通常满足 $\omega \gg \omega_B$ 且 $\omega \gg \omega_p$。我们可以对波数 $k$ 进行泰勒展开：
\begin{align}
    k_{R,L} &= \frac{\omega}{c} \sqrt{1 - \frac{\omega_p^2}{\omega^2(1 \mp \omega_B/\omega)}} \\
            &\approx \frac{\omega}{c} \left[ 1 - \frac{\omega_p^2}{2\omega^2} \left( 1 \pm \frac{\omega_B}{\omega} \right) \right]
\end{align}
计算差值 $(k_R - k_L)$：
\begin{equation}
    k_R - k_L \approx \frac{\omega}{c} \left[ - \frac{\omega_p^2}{2\omega^2} \left( \frac{\omega_B}{\omega} - \left( -\frac{\omega_B}{\omega} \right) \right) \right] = - \frac{\omega_p^2 \omega_B}{c \omega^2}
\end{equation}
将 $\omega_p^2 = 4\pi n e^2/m$ 和 $\omega_B = e B_\parallel / mc$ 代入积分公式（忽略负号，只关心旋转幅度）：
\begin{equation}
    \Delta \theta = \frac{2\pi e^3}{m^2 c^2 \omega^2} \int_0^d n_e B_\parallel ds
\end{equation}

由于 $\omega = 2\pi c / \lambda$，则 $\Delta \theta \propto \lambda^2$。
定义\textbf{旋转量} ：
\begin{equation}
    \text{RM} \equiv \frac{e^3}{2\pi m^2 c^4} \int_0^d n_e B_\parallel ds
\end{equation}
于是观测方程可以写为非常简洁的形式：
\begin{equation}
    \Delta \theta = \text{RM} \cdot \lambda^2
\end{equation}

\begin{derivationnote}[测量方法]
    通过在至少两个不同波长 $\lambda_1, \lambda_2$ 处测量偏振角 $\theta_1, \theta_2$，可以通过斜率求出 RM：
    \[ \text{RM} = \frac{\theta_1 - \theta_2}{\lambda_1^2 - \lambda_2^2} \]
    结合色散量 $\text{DM} = \int n_e ds$，我们可以估算视线方向上的\textbf{平均磁场}：
    \[ \langle B_\parallel \rangle \propto \frac{\text{RM}}{\text{DM}} \]
\end{derivationnote}

\section{高能辐射过程中的等离子体效应}

在讨论高能粒子辐射之前，我们先简要总结入射电磁波（光子）与冷等离子体相互作用的基本效应：
\begin{itemize}
    \item \textbf{反射与传播}: 当入射波频率 $\omega < \omega_p$ 时，波发生全反射；当 $\omega > \omega_p$ 时，波在介质中传播（被散射）。
    \item \textbf{折射}: 若介质非均匀（$n(\vb{r})$ 随位置变化），光线路径会发生弯曲。
    \item \textbf{法拉第旋转}: 线偏振光在磁化冷等离子体中传播时，偏振面会发生旋转。
\end{itemize}
除上述传播效应外，当\textbf{相对论性高能粒子}在介质中运动时，介质的折射特性还会导致两种独特的辐射现象：切伦科夫辐射和 Razin 效应。


\subsection{切伦科夫辐射 (Cherenkov Radiation)}

切伦科夫辐射的本质并非来自粒子的加速度，而是来自介质中麦克斯韦方程组的解在特定速度下的\textbf{奇点性质}。我们可以通过将真空中的李纳-维谢尔（Liénard-Wiechert）势与场直接推广到介质中来理解。
在折射率为 $n=\sqrt{\epsilon}$ 的各向同性均匀介质中（设 $\mu=1$），电磁场的波动方程为：
\begin{equation}
    \left( \nabla^2 - \frac{n^2}{c^2} \pdv[2]{}{t} \right) \phi = - \frac{4\pi \rho}{\epsilon}
\end{equation}
与真空中的波动方程对比，形式完全一致，仅需进行如下\textbf{参数映射}：
\begin{equation}
    c \to c' = \frac{c}{n}, \quad \implies \quad \beta \equiv \frac{v}{c} \to \beta' = \frac{v}{c'} = \beta n
\end{equation}
这意味着，我们可以直接利用真空中的解，通过替换 $\beta \to \beta n$ 来得到介质中的解。


考虑一个以常速度 $\vb{v}$ 运动的带电粒子（加速度 $\dot{\vb{v}} = 0$）。
在真空中，其产生的电场仅包含与距离平方成反比的库仑项（速度场），不包含 $1/R$ 的辐射项。
应用上述代换，介质中的电场 $\vb{E}(\vb{r}, t)$ 表达式变为：
\begin{equation}
    \vb{E}(\vb{r}, t) = \frac{e (1 - \beta^2 n^2)}{ \epsilon R^2 \kappa^3 } \frac{\vb{n} - \vb{v}/(c/n)}{|\vb{n} - \vb{v}/(c/n)|}
    \approx \frac{e (1 - \beta^2 n^2) (\hat{\vb{n}} - \beta n \hat{\vb{v}})}{\epsilon R^2 (1 - \beta n \cos\theta)^3}
\end{equation}
其中关键的\textbf{推迟因子} (Retardation Factor) $\kappa$ 修改为：
\begin{equation}
    \kappa = 1 - \frac{\vb{n} \cdot \vb{v}}{c/n} = 1 - \beta n \cos\theta
\end{equation}

\subsubsection{辐射产生机制}

观察上述电场表达式，我们可以清晰地看到辐射产生的物理判据：

\paragraph{1. 亚光速情况 ($\beta n < 1$)}
分母 $\kappa = 1 - \beta n \cos\theta$ 始终大于 0（因为 $|\cos\theta| \le 1$）。
此时电场 $\vb{E}$ 在所有方向上都是有限的，且随距离以 $1/R^2$ 衰减。这只是一个被介质极化修正了的、跟随粒子移动的库仑场，\textbf{没有辐射}。

\paragraph{2. 超光速情况 ($\beta n > 1$)}
当粒子的速度超过介质中的光速时，分母 $\kappa$ 有可能为零。
存在一个特定的\textbf{切伦科夫角} $\theta_c$，满足：
\begin{equation}
    1 - \beta n \cos\theta_c = 0 \implies \cos\theta_c = \frac{1}{\beta n}
\end{equation}
\begin{itemize}
    \item \textbf{场的奇点}: 在 $\theta = \theta_c$ 的方向上，分母 $\kappa \to 0$，电场强度 $\vb{E} \to \infty$。
    \item \textbf{物理意义}: 这意味着无限多的波前在这一角度上重叠（相位堆积），形成了类似声爆（Sonic Boom）的电磁激波。
    \item \textbf{能量传输}: 尽管粒子没有加速度，但由于场的奇异性，能量不再被束缚在粒子周围，而是沿着 $\theta_c$ 的圆锥面（Mach Cone）向外传播至无穷远，形成切伦科夫辐射。
\end{itemize}


\subsection{Razin 效应 (Razin-Tsytovich Effect)}

Razin 效应是指由于等离子体折射率 $n < 1$，导致相对论性粒子的\textbf{同步辐射在低频段被显著抑制}的现象。

\subsubsection{临界频率的推导}

同步辐射的强度主要源于相对论性的\textbf{集束效应} (Beaming)。在真空中，辐射集中在 $1/\gamma$ 的极小张角内。
在介质中，辐射功率的分母因子为 $(1 - \beta n \cos\theta)$。我们在 $\theta \approx 0$ 附近对其进行泰勒展开。
利用 $\beta \approx 1 - 1/2\gamma^2$，$\cos\theta \approx 1 - \theta^2/2$，以及 $n \approx 1 - \omega_p^2/2\omega^2$：
\begin{align}
    1 - \beta n \cos\theta &\approx 1 - \left(1 - \frac{1}{2\gamma^2}\right) \left(1 - \frac{\omega_p^2}{2\omega^2}\right) \left(1 - \frac{\theta^2}{2}\right) \\
    &\approx \frac{1}{2} \left( \frac{1}{\gamma^2} + \frac{\omega_p^2}{\omega^2} + \theta^2 \right)
\end{align}
\n{忽略了高阶小量。可以看出，介质的存在相当于增加了一个“散焦项” $\omega_p^2/\omega^2$。}

\textbf{截止条件}:
当介质引入的散焦项超过真空中的集束项时，集束效应被破坏，辐射受到抑制。
\begin{equation}
    \frac{\omega_p^2}{\omega^2} > \frac{1}{\gamma^2} \implies \omega < \gamma \omega_p
\end{equation}
因此，$\omega_{Razin} = \gamma \omega_p$ 定义了同步辐射在等离子体中的低频截止频率。
这解释了为什么在致密等离子体源中，观测到的同步辐射谱在低频端会急剧下降，且截断频率远高于普通的等离子体频率 $\omega_p$。